% ========================================================
% BÀI TẬP TOÁN 5 - KẾT NỐI TRI THỨC VỚI CUỘC SỐNG
% Bài 27: Đường tròn. Chu vi và diện tích hình tròn
% Định dạng: BTTeX Part III - Trắc nghiệm điền số ngắn
% ========================================================

% Câu 1
\begin{ex}
  Đường kính của hình tròn gấp mấy lần bán kính?
  \shortans{2}
  \loigiai{
    $d = 2 \times r$. Đường kính gấp 2 lần bán kính.
  }
\end{ex}

% Câu 2
\begin{ex}
  Bán kính hình tròn là $5\text{ cm}$. Đường kính của hình tròn đó là bao nhiêu xăng-ti-mét?
  \shortans{10}
  \loigiai{
    Đường kính là $5 \times 2 = 10\text{ cm}$.
  }
\end{ex}

% Câu 3
\begin{ex}
  Đường kính hình tròn là $16\text{ dm}$. Bán kính của hình tròn đó là bao nhiêu đề-xi-mét?
  \shortans{8}
  \loigiai{
    Bán kính là $16 : 2 = 8\text{ dm}$.
  }
\end{ex}

% Câu 4
\begin{ex}
  Một hình tròn có bán kính $10\text{ cm}$. Chu vi của hình tròn đó là bao nhiêu phần mười xăng-ti-mét? (Lấy $\pi = 3{,}14$, tính chu vi nhân 10).
  \shortans{628}
  \loigiai{
    Chu vi là $2 \times 3{,}14 \times 10 = 62{,}8\text{ cm}$. Nhân 10 được 628.
  }
\end{ex}

% Câu 5
\begin{ex}
  Một hình tròn có đường kính $20\text{ cm}$. Chu vi hình tròn đó là $62{,}8\text{ cm}$. Phần nguyên của chu vi là bao nhiêu?
  \shortans{62}
  \loigiai{
    Phần nguyên của $62{,}8$ là 62.
  }
\end{ex}

% Câu 6
\begin{ex}
  Một hình tròn có bán kính $10\text{ cm}$. Diện tích của hình tròn đó là bao nhiêu xăng-ti-mét vuông? (Lấy $\pi = 3{,}14$).
  \shortans{314}
  \loigiai{
    $S = 3{,}14 \times 10 \times 10 = 314\text{ cm}^2$.
  }
\end{ex}

% Câu 7
\begin{ex}
  Một hình tròn có bán kính $20\text{ dm}$. Diện tích của hình tròn đó là bao nhiêu đề-xi-mét vuông? (Lấy $\pi = 3{,}14$).
  \shortans{1256}
  \loigiai{
    $S = 3{,}14 \times 20 \times 20 = 1\,256\text{ dm}^2$.
  }
\end{ex}

% Câu 8
\begin{ex}
  Một hình tròn có đường kính $10\text{ m}$. Diện tích hình tròn đó nhân với 4 là bao nhiêu mét vuông? (Lấy $\pi = 3{,}14$).
  \shortans{314}
  \loigiai{
    Bán kính là 5 m. $S = 3{,}14 \times 25 = 78{,}5\text{ m}^2$. $78{,}5 \times 4 = 314\text{ m}^2$.
  }
\end{ex}

% Câu 9
\begin{ex}
  Nếu bán kính hình tròn tăng lên gấp đôi thì chu vi của nó tăng lên mấy lần?
  \shortans{2}
  \loigiai{
    Chu vi tỉ lệ thuận với bán kính nên tăng 2 lần.
  }
\end{ex}

% Câu 10
\begin{ex}
  Nếu bán kính hình tròn tăng lên gấp đôi thì diện tích của nó tăng lên mấy lần?
  \shortans{4}
  \loigiai{
    Diện tích tỉ lệ với bình phương bán kính: $2 \times 2 = 4$ lần.
  }
\end{ex}

% Câu 11
\begin{ex}
  Nếu bán kính hình tròn tăng lên gấp 3 lần thì diện tích tăng lên mấy lần?
  \shortans{9}
  \loigiai{
    $3 \times 3 = 9$ lần.
  }
\end{ex}

% Câu 12
\begin{ex}
  Một bánh xe đạp có bán kính $0{,}35\text{ m}$. Đường kính bánh xe đạp đó là bao nhiêu phần mười mét? (Điền đường kính nhân 10).
  \shortans{7}
  \loigiai{
    Đường kính là $0{,}35 \times 2 = 0{,}7\text{ m}$. Nhân 10 được 7.
  }
\end{ex}

% Câu 13
\begin{ex}
  Một bánh xe đạp có chu vi $2{,}2\text{ m}$. Khi bánh xe lăn được $100$ vòng trên đường thì xe đi được quãng đường bao nhiêu mét?
  \shortans{220}
  \loigiai{
    Quãng đường là $2{,}2 \times 100 = 220\text{ m}$.
  }
\end{ex}

% Câu 14
\begin{ex}
  Một bánh xe ô tô lăn được $500$ vòng thì đi được quãng đường $1\,000\text{ m}$. Chu vi bánh xe ô tô đó là bao nhiêu mét?
  \shortans{2}
  \loigiai{
    Chu vi là $1\,000 : 500 = 2\text{ m}$.
  }
\end{ex}

% Câu 15
\begin{ex}
  Một hồ nước hình tròn có bán kính $10\text{ m}$. Diện tích hồ nước là bao nhiêu mét vuông? (Lấy $\pi = 3{,}14$).
  \shortans{314}
  \loigiai{
    Diện tích là $3{,}14 \times 10 \times 10 = 314\text{ m}^2$.
  }
\end{ex}

% Câu 16
\begin{ex}
  Một bồn hoa hình tròn có đường kính $4\text{ m}$. Diện tích bồn hoa nhân với 100 là bao nhiêu? (Lấy $\pi = 3{,}14$).
  \shortans{1256}
  \loigiai{
    Bán kính 2 m. $S = 3{,}14 \times 4 = 12{,}56\text{ m}^2$. Nhân 100 được 1256.
  }
\end{ex}

% Câu 17
\begin{ex}
  Một mặt đồng hồ hình tròn có bán kính $15\text{ cm}$. Chu vi mặt đồng hồ là $94{,}2\text{ cm}$. Chữ số hàng phần mười là bao nhiêu?
  \shortans{2}
  \loigiai{
    Chữ số hàng phần mười là 2.
  }
\end{ex}

% Câu 18
\begin{ex}
  Một hình tròn có chu vi là $31{,}4\text{ cm}$. Đường kính hình tròn đó là bao nhiêu xăng-ti-mét? (Lấy $\pi = 3{,}14$).
  \shortans{10}
  \loigiai{
    Đường kính là $31{,}4 : 3{,}14 = 10\text{ cm}$.
  }
\end{ex}

% Câu 19
\begin{ex}
  Ở câu trên, bán kính của hình tròn là bao nhiêu xăng-ti-mét?
  \shortans{5}
  \loigiai{
    Bán kính là $10 : 2 = 5\text{ cm}$.
  }
\end{ex}

% Câu 20
\begin{ex}
  Một hình tròn có diện tích $314\text{ cm}^2$. Bán kính của hình tròn đó là bao nhiêu xăng-ti-mét? (Lấy $\pi = 3{,}14$).
  \shortans{10}
  \loigiai{
    $r \times r = 314 : 3{,}14 = 100 \implies r = 10\text{ cm}$.
  }
\end{ex}

% Câu 21
\begin{ex}
  Một hình tròn có diện tích $78{,}5\text{ cm}^2$. Bán kính hình tròn đó là bao nhiêu xăng-ti-mét? (Lấy $\pi = 3{,}14$).
  \shortans{5}
  \loigiai{
    $r \times r = 78{,}5 : 3{,}14 = 25 \implies r = 5\text{ cm}$.
  }
\end{ex}

% Câu 22
\begin{ex}
  Hai hình tròn có bán kính lần lượt là $2\text{ cm}$ và $4\text{ cm}$. Tỉ số diện tích của hình tròn lớn so với hình tròn bé là bao nhiêu?
  \shortans{4}
  \loigiai{
    Tỉ số diện tích là $(4 : 2)^2 = 2^2 = 4$.
  }
\end{ex}

% Câu 23
\begin{ex}
  Một cái nong hình tròn có đường kính $1\text{ m}$. Bán kính cái nong là bao nhiêu đề-xi-mét?
  \shortans{5}
  \loigiai{
    $1\text{ m} = 10\text{ dm}$. Bán kính là $10 : 2 = 5\text{ dm}$.
  }
\end{ex}

% Câu 24
\begin{ex}
  Nửa hình tròn có bán kính $10\text{ cm}$. Diện tích của nửa hình tròn đó là bao nhiêu xăng-ti-mét vuông? (Lấy $\pi = 3{,}14$).
  \shortans{157}
  \loigiai{
    Diện tích cả hình tròn là 314. Nửa hình tròn là $314 : 2 = 157\text{ cm}^2$.
  }
\end{ex}

% Câu 25
\begin{ex}
  Bán kính hình tròn là $100\text{ m}$. Diện tích hình tròn đó là bao nhiêu héc-ta? (Lấy $\pi = 3{,}14$).
  \shortans{3}
  \loigiai{
    Làm tròn đến hàng đơn vị: $S = 3{,}14 \times 100 \times 100 = 31\,400\text{ m}^2 = 3{,}14\text{ ha}$. Phần nguyên là 3.
  }
\end{ex}
